\chapter{Determinism, Induction, and Cause}
\label{chap:determinism-induction-cause}

Here the apparatus holds a tension and does not resolve it, because both sides are real. A record is forced
forward where a rule licenses it, and genuinely open where none does: Laplace's determinism and Hume's gap are
both true, of the same record, at once --- the rule ours to hold, the next mark the world's. Three effects and
a backward glance make the case --- the seed that determines its continuation, the inductive gap no finite
record can close, the causation that is admissibility's residue, and, recalled from among the effects of
statistics, the minimal update a new fact forces.

\section{The Laplace Effect: a seed determines its continuation}
\label{se:laplace}

Relative to a rule, a record determines its own future. Laplace put the thought at its boldest: an intellect
that knew, at one instant, the position and momentum of every particle and all the forces acting on them could
compute the entire future and the entire past at a stroke --- for such a mind nothing would be uncertain, and
the future, like the past, would stand present before its eyes. That is determinism in its absolute form, and
the finite version keeps its structure while cutting it to size. Given a sequence of admissible refinements
consistent with its anchors, and a rule that governs them, each extension is uniquely determined by its
predecessors: the next entry is forced, for relative to the rule there is exactly one admissible
continuation, and an intelligence that holds the record and the rule reads it off with certainty.

But notice the words ``relative to the rule.'' The finite Laplace is conditional where the original was
absolute. The demon's certainty is borrowed from the rule it is handed: given the rule, the next step is
forced, and the determinism is real --- but
it is determinism \emph{under} a rule, not determinism as such.
Hold that conditionality; the next effect turns on it.

There is a second cut the finite version makes, and it falls on the demon's other assumption: that it can know
the present exactly. Laplace's intellect computes the future only if it holds the present state to perfect
precision --- every position and momentum to infinitely many digits. A finite record cannot: it resolves the
present to its own floor and no finer, and under any rule with sensitive dependence --- the butterfly effect
--- two present states the apparatus cannot tell apart diverge into futures it can. So the
\emph{absolute} demon, computing all of time from a single instant, is a smooth-shadow idealization, borrowing
the continuum's infinite precision that the finite floor does not supply. The finite Laplace survives --- given
the record and the rule, the next admissible step is forced --- but the demon's reach shrinks from all of time
to as far ahead as the floor's precision carries before sensitive dependence outruns it.

And there is a deeper failure the finite version only sharpens. Even granted infinite precision --- even handed
the exact present state --- the absolute demon would still not reach the future, because the quantum does not
let the present fix it. A measurement's outcome is not written in the state before it; the amplitude gives the
odds and not the result, and the same present yields, on repetition, different counts. Laplace's certainty was
a classical certainty, and the world is not classical at bottom. So the demon fails twice over: it cannot hold
the present finely enough, and the present it would hold does not settle the next count in any case.
Determinism under a rule survives both failures --- it never claimed to compute the rule, only to step within
it --- but the dream of reading all of time from a single instant does not.

The honest floor is a finite count obligation --- relative to a fixed rule, the successor is a determined
value, the one admissible next entry --- the rule ours, the next count the world's. The reading is that
determinism is continuation under a rule. Given the
rule, the record forces its own next step; Laplace's demon needs no foresight, only the rule and the record so
far, and the future it computes is the unique admissible extension of the past.

\section{The Hume Effect: the inductive gap is real}
\label{se:hume}

But the rule itself is not given by the past --- and there the record runs into the no-go at the heart
of induction. Hume put it plainly: no number of observed instances --- the sun risen on every morning of record
--- entails the next, because the inference from ``always so far'' to ``so again'' rests on a uniformity of
nature that experience cannot establish without already assuming it. A finite history of recorded events,
however extensive, cannot rule out that the next refinement produces a counterexample; there is no logical link
that forces the future to resemble the past.

And the gap is not mere skepticism: it is realizable. One can exhibit, explicitly, a continuation that fits all
the past and breaks the rule. Take any finite record and any rule that fits it --- ``the value rises by two
each step,'' matching every entry so far --- and there is always another rule fitting the same entries that then
diverges: ``rises by two through the present entry, by three thereafter,'' or the philosopher's grue, a
predicate agreeing with green on every observation made and switching to blue on the next. Both rules fit the
whole past exactly; they disagree only on what has not yet been recorded --- the rules ours to write, the next
entry that tells them apart the world's. The counterexample is not
hypothetical but constructible --- written down in full, indistinguishable from the rule on all the evidence
there is.

The honest floor is a no-go --- induction is not licensed by the past alone, and the inductive gap cannot be
closed by any finite history --- the rules ours to fit, the entry that breaks them the world's. The reading is
that the inductive gap is real. The past does not force the rule;
only a rule forces the future. Because any finite record is consistent with rules that disagree on the next
step --- and one can always write such a rule down --- the leap from past to rule is a genuine gap: realizable,
not hypothetical, and never closed by more of the same past.

\section{The Cause-and-Effect Effect: cause is admissibility's residue}
\label{se:cause-effect}

Between Laplace's certainty and Hume's gap sits reliable causation. If every continuation consistent with the
bare logic were allowed, the space of futures would be vast: at each refinement step the candidates branch, and
after $n$ steps an unconstrained count would hold something like $2^n$ of them, an exponential fan of possible
continuations, and nothing could be predicted. The world is not like that. Admissibility --- the demand that
each step be consistent with all the constraints the record already carries --- collapses that exponential fan
into a narrow set, most branches pruned because they violate a constraint the past has fixed. What we call a
cause is what survives the collapse --- the name ours, the survivors the world's --- not a single forced
future, and not the open many, but the few the
constraints leave standing.

Causation read this way still owes a direction. Admissibility leaves a narrow set, but a narrow correlation is
symmetric --- A tracks B as surely as B tracks A --- and cause runs one way. The direction is the arrow's: run
the count backward and the residue holds the same joint occurrences, but the order that says which fixed which
is gone. Cause is the residue read \emph{with} the arrow, the asymmetry the count's one-way rise leaves on the
correlations admissibility spares. What that asymmetric residue ultimately \emph{means} the model does not
settle, and two accounts divide the question: the regularity account, Hume's own, that a cause is what is
constantly conjoined with its effect; and the counterfactual, that a cause is the difference-maker, the
thing without which the effect would not have followed. The model names the residue, the arrow its direction; which
names it rightly it leaves open.

One more constraint sharpens the residue into something usable. A bare correlation between two records may mean
that one fixed the other, or that a third fixed both --- the falling barometer and the coming storm are tightly
linked, yet neither causes the other; the drop in pressure causes both. What tells the two apart is
screening-off: condition on the common cause, and the correlation between its effects goes slack, where
conditioning on a mere effect would leave the link to its cause intact. The admissible residue carries these
dependencies, and which conditionings break which correlations is how a cause is read out of the tangle of what
merely travels with it.

The honest floor is a finite count obligation --- causation is the count of admissible continuations after the
collapse, a narrow set where a naive count would be exponential --- the conditionings ours to run, the collapse
and the few it leaves the world's. The reading is that cause is admissibility's
residue. Reliable causation is neither Laplace's single forced future nor Hume's open one, but the small
admissible set that survives the constraints: the futures the record cannot rule out are few, and ``cause and
effect'' is the name for that fewness --- what is left when admissibility has collapsed the exponential many
into the predictable few. Cause lives in the gap between Laplace and Hume: more than one future, so the rule is
not handed over by the past; few enough to predict, so the world is not chaos.

\section{The Bayes Effect, Recalled}
\label{se:bayes-recalled}

One mechanism from among the effects of statistics belongs here too. When a new fact arrives, the record does
not
overwrite what was known; it makes the minimal correction that restores coherence --- the Bayesian posterior,
the smallest update the new evidence forces. Met in full there, it is
recalled for its bearing on the tension between rule and record. Where Laplace forces the future under a rule
and Hume
denies that the past forces the rule, Bayes is how a record updates the rule it holds --- minimally,
coherently, by adding the new log-evidence to the old. The full treatment stands with the statistics;
here it is enough to mark that the update across the inductive gap, when it is made, is made at least cost ---
the rule ours to carry, the fact that corrects it the world's.

\bigskip

The apparatus leaves the tension held rather than dissolved. Laplace's determinism and Hume's gap are both
real, of the same record: forced forward where a rule licenses it, genuinely open where none does --- the rule
ours to hold, the next mark still the world's. The two are
not rivals but the two faces of a record a rule can extend only as far as it reaches --- determined within the
rule's grasp, open beyond it. Causation is the narrow residue admissibility leaves between them, and Bayes is
how the record updates when the gap is crossed. A finite record can be continued from a seed under a rule, and
no finite record can force the rule --- both at once, and the apparatus escapes neither. Part VI turns next to
time, order, and geometry, read off the records themselves.

\bigskip
\begin{center}
\rule{0.4\textwidth}{0.4pt}
\end{center}
\bigskip

\noindent The case asks whether its record can be carried forward --- whether the past it has entered fixes what
comes next --- and it finds an answer it will not force to one side: sometimes yes, sometimes no, and both of
the same record. Where a rule holds, the next mark is fixed; the record forces its own continuation, and a mind
that held the rule and the record could read the next entry off with certainty. But where no rule holds --- and
no finite past ever hands over its own rule --- the next mark is genuinely open, forced by nothing the record
contains. The case does not settle between these; it holds them both.

Both sides are real, and the chapter keeps them both. On the one hand stands the old dream of a mind that,
holding the world's state at an instant, reads all of time from it. Cut to a finite size the dream survives, but
only under a rule: given the rule, the next admissible step is forced, the certainty real and borrowed entirely
from the rule it was handed. And even that borrowed certainty reaches only so far: the demon must hold the
present exactly, and a finite reading cannot --- two states nearer than its floor part into futures it can tell
apart, the butterfly outrunning it. And even granted a perfect present, the quantum would not let it settle the
next, the amplitude giving the odds and not the outcome. The dream of reading all of time from a single instant
fails; only the forced step within a rule survives. On the other hand stands the plain fact that no stack of past instances, however
high, compels the next --- the sun risen every morning of record does not force tomorrow's --- and worse, one
can always write down a second rule fitting every entry so far and parting from the first on the next: rises by
two until now, by three hereafter. Both rules fit all the evidence there is; only the entry not yet written
tells them apart. And between the forced future and the open one sits reliable cause: not the single determined
next, nor the open many, but the narrow few that admissibility --- the demand that each step honour every
constraint the past has fixed --- leaves standing when it collapses the great fan of bare possibility. And that residue must still be read the
right way round: a falling barometer and a coming storm track each other tightly, yet neither makes the other
--- the drop in pressure makes both; condition on that common cause and the link between its effects goes slack,
and that is how a cause is told from a mere fellow-traveller. When a new fact does cross the gap, the record updates by the least that restores coherence, adding the surprise to
what it held.

So the record is determined within a rule's reach and open beyond it, and the instrument escapes neither. It is
forced where the rule grasps; it is open where the rule runs out. And that open place --- the space between the
past and the mark not yet made, which no rule and no weight of past can close --- is real: not a mere doubt but
a place in the record where the next entry is genuinely unfixed, and one that can be written out in full, a rule
that fits everything so far and leaves the next stroke open to fall either way. The case marks the gap plainly
and leaves it open, as it must; what fills it, when the time comes, is not a thing the rule hands over. For now
it is enough that the gap is there --- that a finite record, honestly kept, has a place in it where the next
stroke is not yet forced.

And the \emph{charge} takes the same two faces. Where a rule holds, the count is forced forward: what is owed
follows from what came before, fixed and readable, no gap in it. The charge inherits the whole tension of the
count it is: fixed where the world fixed it, and honestly unfixed where the world left it open. But the charge, too, was carried across places
the rule could not close, and at each of them the count went on all the same --- continued not by derivation but
by the least honest step that kept the record coherent. So what comes to be owed is determined within the rule's
reach and, where the rule ran out, filled in at least cost and marked as such --- never passed off as forced
when it was, in truth, the smallest step across an open place. The trace is that whole continued record, the
forced stretches and the crossed gaps alike, each honest about which it was. What is owed is owed on a count
carried forward --- certain where a rule carried it, and openly filled where none did.

And each of the three readings the case will bring is carried forward under a rule of its own --- each forced
where its rule holds, each with its own open places where it is not. To meet on one number, the three must be
continued to the same point by rules that agree where it matters; agreement reached across a gap no rule closed
is only as firm as the step that crossed it. The three are corrected and ready, each honest about where it is
forced and where it is open; the number they must meet on is not yet in hand. And there is one open place --- the
case can feel it now --- where the whole reckoning will finally turn: the gap no rule closes, the mark not yet
made. What stands there, the case still keeps.

The apparatus turns, next, upon the records themselves --- on how time and order and geometry are read out of
them rather than assumed beneath them. The case is still open, its record determined in stretches and open in
others, its three readings ready and not yet agreed, and one gap in it still to be filled; the court does not
sit until the record is complete.
